\section{Time-Sliced Transaction--Address Knowledge Graph}\label{sec:kg}

\subsection{Schema}
The TS-KG (Fig.~\ref{fig:schema}) has two entity types. A \emph{Transaction}
node exists for each Elliptic transaction. An \emph{Address@t} node exists for
each pair of an address and a time step in which it is active (920{,}691
instances of 822{,}942 addresses). There are three relation types:
\texttt{FLOWS\_TO} between transactions, \texttt{SPENDS\_IN} from an input
address to a transaction, and \texttt{PAYS} from a transaction to an output
address. Instantiating addresses per step keeps every edge inside one step.
The training, validation and test periods are therefore disconnected
sub-graphs, so message passing can never carry information from the test
period into training (strictly inductive evaluation), and no address
statistic computed at a later date ever enters the graph.

An address's \emph{history} enters only through causal attributes of its
Address@t node: whether and how often it was seen in earlier steps, steps
since first and last appearance, cumulative earlier input and output counts,
and its input and output counts within step $t$. Within step $t$ we treat the
transactions of the step as observable, as in the original benchmark; this
corresponds to scoring transactions at the end of each step, a fixed-delay
regime that~\cite{safar2026} identifies as legitimate.

\begin{figure}[t]
\centering
\begin{tikzpicture}[
  tx/.style={rectangle, rounded corners=2pt, draw=blue!60!black, fill=blue!8, minimum width=9mm, minimum height=5mm, font=\scriptsize},
  ad/.style={ellipse, draw=orange!70!black, fill=orange!12, minimum width=10mm, minimum height=5mm, inner sep=1pt, font=\scriptsize},
  rel/.style={-{Latex[length=1.6mm]}, thick},
  lab/.style={font=\tiny, fill=white, inner sep=1pt}]
\node[tx] (t1) at (0,0) {Tx$_1$};
\node[tx] (t2) at (2.6,0) {Tx$_2$};
\node[tx] (t3) at (5.2,0) {Tx$_3$};
\node[ad] (a1) at (1.3,1.35) {A@$t$};
\node[ad] (a2) at (3.9,1.35) {B@$t$};
\node[ad] (a3) at (2.6,-1.35) {C@$t$};
\draw[rel, gray] (t1) -- (t2) node[lab, midway, below=1pt] {FLOWS\_TO};
\draw[rel, orange!70!black] (a1) -- (t1) node[lab, midway, left=1pt] {SPENDS\_IN};
\draw[rel, orange!70!black] (a1) -- (t2);
\draw[rel, blue!60!black] (t2) -- (a2) node[lab, midway, left=1pt] {PAYS};
\draw[rel, orange!70!black] (a2) -- (t3);
\draw[rel, blue!60!black] (t1) -- (a3);
\draw[rel, blue!60!black] (t3) -- (a3);
\node[draw=gray, dashed, rounded corners, fit=(t1)(t3)(a1)(a2)(a3), inner sep=5pt] (box) {};
\node[font=\tiny, anchor=north west, gray] at (box.north west) {time step $t$};
\end{tikzpicture}
\caption{Time-sliced knowledge graph schema. Address nodes are instantiated per
time step; history enters only through causal node attributes. Two-hop paths
through an address define meta-relations: Tx$_1$ and Tx$_2$ are
\texttt{CO\_SPEND} (shared input A), Tx$_1$ and Tx$_3$ are
\texttt{CO\_RECEIVE} (shared output C), Tx$_2$ and Tx$_3$ are \texttt{CHAIN}
(Tx$_2$ pays B, which spends in Tx$_3$).}
\label{fig:schema}
\end{figure}

\subsection{Meta-relations and homophily}
Two-hop paths through an Address@t node define three meta-relations between
transactions: \texttt{CO\_SPEND} (a shared input address),
\texttt{CO\_RECEIVE} (a shared output address) and \texttt{CHAIN} (one
transaction pays an address that another spends from). Address instances
linked to more than 200 transactions in a step (67 exchange-like hubs) are
excluded from the meta-relations. Table~\ref{tab:homophily} reports, for
labelled transactions, the class-adjusted homophily~\cite{platonov2023}
(0 for random wiring, 1 for perfect homophily), together with the share of
illicit transactions that have at least one labelled neighbour (coverage) and
the share of their labelled neighbours that are illicit.

\begin{table}[t]
\caption{Label homophily of transaction relations (labelled nodes).}
\label{tab:homophily}
\centering
\footnotesize
\setlength{\tabcolsep}{3.2pt}
\resizebox{\columnwidth}{!}{%
\begin{tabular}{@{}lrrrrrr@{}}
\toprule
& \multicolumn{3}{c}{Train (steps 1--34)} & \multicolumn{3}{c}{Test (steps 35--49)} \\
\cmidrule(lr){2-4}\cmidrule(l){5-7}
Relation & Adj.\ hom. & Cov.\ & Ill.\ nbr & Adj.\ hom. & Cov.\ & Ill.\ nbr \\
\midrule
FLOWS\_TO   & 0.553 & 0.551 & 0.581 & 0.376 & 0.532 & 0.395 \\
CO\_SPEND$^{a}$ & 1.000 & 0.186 & 1.000 & 1.000 & 0.125 & 1.000 \\
CO\_RECEIVE & 1.000 & 0.759 & 1.000 & 1.000 & 0.667 & 1.000 \\
CHAIN       & 0.865 & 0.527 & 0.871 & 0.309 & 0.547 & 0.315 \\
KG union    & \textbf{0.954} & \textbf{0.906} & \textbf{0.957} & \textbf{0.915} & \textbf{0.930} & \textbf{0.921} \\
\bottomrule
\multicolumn{7}{@{}l}{\scriptsize Cov.: illicit transactions with a labelled neighbour. Ill.\ nbr: illicit share of}\\
\multicolumn{7}{@{}l}{\scriptsize their labelled neighbours. $^{a}$Homophilous by construction (labels follow inputs).}
\end{tabular}}
\end{table}

Three observations motivate graph learning on the TS-KG. First, the
transaction graph used by prior GNN work is only moderately homophilous and
becomes much less so after the dark-market shutdown in the test period
(adjusted homophily 0.55$\rightarrow$0.38), consistent with its weak results
under strict evaluation~\cite{maganti2026}. Second, the KG relations connect
93\% of illicit test transactions to labelled neighbours, of which 92\% are
illicit; this is the regime in which message passing is expected to
help~\cite{zhu2020,platonov2023}. Third, part of this signal follows directly
from how labels were made: transactions sharing an input address share the
labelling entity, so \texttt{CO\_SPEND} is perfectly homophilous by
construction. It covers only 13\% of illicit test transactions, however,
whereas \texttt{CO\_RECEIVE}, which is not determined by the labelling rule,
covers 67\%. In practice the shared output addresses are change addresses and
deposit addresses of the same entity, the structure that forensic address
clustering exploits~\cite{meiklejohn2013}. We return to this caveat in
Section~\ref{sec:discussion}.

\textbf{Address reuse across time is a licit signal.} Only 17 of the 1{,}083
illicit test transactions touch an address that was active during training,
against 5{,}833 of the 16{,}670 labelled test transactions overall. Illicit
actors use fresh addresses, so cross-period address knowledge mainly
identifies licit activity. Consequently, the \emph{unseen-entity} test view
recommended by~\cite{safar2026} (test transactions none of whose addresses
appeared in training) retains 98\% of the illicit test transactions.
